%  The technical layer of the two-layer split.
%
%  The article states what each stage does and why; this states the arithmetic
%  it does it with -- the cost function, every threshold, the round schedule,
%  and the solver detail, for the 2-D reconstruction and the depth stage. The
%  division is by *depth of detail*, not by importance: nothing here is
%  optional to the method, and every number here is the one the released
%  configuration actually carries.
%
%  Included by supplementary.tex ahead of the measure definitions, so a reader
%  who wants the mechanism meets it before the scoring.

\section{Reconstruction: cost function, gates and schedule}
\label{supp:recon}

This section gives the arithmetic behind Section~2.4 of the article. Values are
those of the released configuration; nothing here was tuned on held-out or
locked data.

\subsection{Skeleton graph}

The mask $M$ is binarised at $>0$, skeletonised, and spurs no longer than
\PlectaParamSpurPx{}\,px are removed iteratively. Skeleton pixels of degree at
least three form junction clusters, and adjacent clusters merge into nodes:
connectors up to \PlectaParamAbsorbPx{}\,px are absorbed into one node,
connectors up to \PlectaParamJoinPx{}\,px merge nearby nodes while remaining
real arms, and free debris up to \PlectaParamFreeDebrisPx{}\,px is absorbed.
Each remaining connected path is stored as an ordered arm carrying one stub at
each end.

\subsection{Stub frames}

A stub frame holds its tip, outward tangent and signed curvature, estimated by
a least-squares arclength polynomial fit read outward from the tip. The fit is
quadratic when the span reaches \PlectaParamMinQuadratic{}\,px and linear
otherwise, and a frame counts as reliable only with at least three samples
spanning at least \PlectaParamSpurPx{}\,px. Local frames use
\PlectaParamWindowLocal{}\,px of support; once links form chains that support
expands to \PlectaParamWindowChain{}\,px, so a stub too short to have a
direction of its own borrows one from the chain it has joined.

\subsection{Weights and scales of the pairing cost}

The cost function itself is stated in the article, which also gives the
reasoning behind its three deliberate properties. What follows are the values
it is evaluated with.

Junction and gap links use separate weight sets, because a junction link is
scored on direction alone while a gap link must also answer for the distance it
spans. The curvature term carries a constant of \PlectaParamKappaScale{}, which
places curvature, in px$^{-1}$, on the scale of the angular terms; without it a
term in reciprocal pixels and a term in radians could not be added. Junction
links set
$w_\ell=0$, a junction having no gap to pay for, and $\ell_0$ and $d_0$ set the
length scale of the gap penalty and the width of the chord fade respectively.

\begin{table}[tbp]
  \centering
  \caption{Cost weights and scales, read from the frozen configuration that
  produced every published number. Junction and gap links are scored with
  separate sets; the gap set alone carries a length term, a junction having no
  gap to pay for.}
  \label{tab:cost-weights}
  \begin{tabular}{llcc}
\toprule
Symbol & Term & Junction & Gap \\
\midrule
$w_d$ & reversal angle $\theta$ & 0.6 & 0.2 \\
$w_t$ & chord turns $\varphi_a+\varphi_b$ & 1.2 & 1 \\
$w_\kappa$ & curvature continuity & 1 & 0.5 \\
$w_\ell$ & length $d/\ell_0$ & --- & 0.15 \\
$\ell_0$ & length scale (px) & --- & 60 \\
$d_0$ & chord-fade scale (px) & 4 & 3 \\
$p_i$ & price of leaving a stub free & 0.62 & 0.3 \\
\bottomrule
\end{tabular}

\end{table}

\subsection{Prices, gates and admissibility}

Leaving a stub unmatched is priced rather than forbidden, and the price $p_i$
also fixes admissibility: an edge is offered only when $C_{ij}<p_i+p_j$. Exact
maximum-weight partial matching then maximises $\sum(p_i+p_j-C_{ij})$ over
disjoint pairs, which minimises total link cost plus $p_i$ for every stub left
free.

The admissibility limit is not a separate threshold. An edge with
$C_{ij}\ge p_i+p_j$ contributes nothing positive to the objective, so a
maximum-weight \emph{partial} matching would never select it, and declining to
offer it removes an edge that could not have been chosen.

At the configured prices junction stubs carry
$p_i=\PlectaParamJunctionPrice{}$, so a junction edge needs
$C<\PlectaParamJunctionLimit{}$. Gap stubs carry
$p_i=\PlectaParamGapPrice{}$, so a gap edge needs
$C<\PlectaParamGapLimit{}$, and must in addition satisfy
$d\le\PlectaParamGapMaxLen{}$\,px, $\theta\le\PlectaParamGapMaxTheta{}$\,rad
and $\varphi\le\PlectaParamGapMaxPhi{}$\,rad at each end. The gap gates are the
tighter set because a gap link asserts a continuation across evidence that is
missing rather than merely ambiguous. The values were fixed on development
scenes.

\subsection{The round schedule}

Rounds $r=0,\ldots,\PlectaParamRounds{}-1$ apply a schedule fraction
$\alpha_r=\PlectaParamAnnealStart{}+(1-\PlectaParamAnnealStart{})\,r/7$ that
scales the unmatched prices and the two gap angle limits, and nothing else.
Scaling a price \emph{down} admits fewer edges, admissibility being
$C_{ij}<p_i+p_j$, and scaling an angle limit down narrows the gap gates, so
early rounds are the conservative ones and are deliberately under-linked: while
the frames are least reliable, only pairings that survive the reduced prices and
narrowed gates are committed, and the constraints relax towards their configured
values as the chains assembled so far sharpen the frames. The final round,
$\alpha_7=1$, is the only one run at the configured parameters, and is the one
returned.

In full:
\begin{enumerate}[nosep,leftmargin=1.5em,label=\arabic*.]
  \item Skeletonise $M$, prune spurs, merge junction clusters into nodes, and
        store the remaining paths as arms with a stub at each end.
  \item For $r=0,\ldots,7$, re-solving both matchings from scratch and carrying
        only the chain structure forward from round $r-1$:
  \begin{enumerate}[nosep,leftmargin=1.5em,label=(\alph*)]
    \item Re-fit every stub frame along the chain it now belongs to.
    \item At each node, form candidate edges between stubs of different arms
          with finite frames and cost, keep those admissible at $\alpha_r$, and
          solve one maximum-weight matching per node.
    \item From the still-unmatched reliable stubs, form gap candidates on
          different arms and different nodes that pass the distance gate and
          the $\alpha_r$-scaled angle gates, and solve one global
          maximum-weight matching over the admissible ones.
  \end{enumerate}
  \item Return the links of round $r=7$.
\end{enumerate}

\subsection{Open-curve guard}

Instances are assumed to be open projected curves, and a guard enforces it:
after each of the two matchings, any matched component whose links close a ring
has its longest link removed. On the \PlectaDevelopmentSceneCount{} development
scenes the guard never fired --- zero links removed --- so it protects against
a configuration the data did not produce rather than shaping the result. Closed
or looping filaments are outside the scope of this paper.

\subsection{Why the junction is solved as a whole}

Joint and sequential solutions coincide at plain crossings and diverge only as
nodes grow dense. Over 20 development scenes they never disagreed at nodes of
degree 2--4 ($n=542$), disagreed at 22\,\% of nodes with 5--8 stubs ($n=187$)
and at 69\,\% of nodes with 9 or more ($n=233$), where the joint solution was
better by a median of 0.90 in total cost.

That comparison holds one node's costs fixed. Run inside the rounds, where each
matching sets the frames the next is scored against, a cheap early choice
propagates and the difference is no longer confined to dense nodes.

\subsection{Instance output}

Connected components of accepted arm links define instances. PLECTA paints
every arm pixel and every junction cluster touched by a chain into that
instance, so intersecting instances share crossing pixels. An isolated one-arm
component shorter than 6\,px is omitted. Accepted gap links determine identity
but are not painted into the fixed centreline output. All effective parameters,
including values previously inherited as code defaults, are materialised in the
released JSON configuration.

\section{Depth stage: evidence and solver}
\label{supp:depthstage}

This section gives the detail behind Section~2.6 of the article.

\subsection{Sampling geometry at a crossing}

Where two instance centrelines intersect, a short arc of each is taken through
the crossing --- the shared \emph{core}, whose half-size follows the two radii
so that it spans both silhouettes --- together with the \emph{flanks} on either
side of it. Flank samples are kept clear of the other filament's stroke,
because a sample taken inside it measures the union of two rods rather than one
of them. A crossing with too few usable flank samples abstains rather than
being decided on thin evidence.

\subsection{The two channels}

Two quantities are read on core and flanks alike: median intensity, and median
gradient magnitude sampled on the filament's own two edges, the axis gradient
being near zero and therefore uninformative. Sharpness continuity is tested per
rod on that rod's own geometry: the core-adjacent samples of rod $k$ sit on
$k$'s edges when $k$ is on top, and inside the other rod's stroke when $k$ is
underneath, where its edge is overwritten.

The filament whose flank appearance the core matches better is called upper,
and the normalised difference between the two channels passes through a
logistic to give $p$. Where the two filaments genuinely look alike the
difference is near zero, $p$ sits near $0.5$, and the crossing abstains; the
width of that abstention band is a configured parameter, as are the two channel
weights.

\subsection{One order for the whole field}

Each decided crossing contributes a directed edge weighted by
$\lvert\operatorname{logit}p\rvert$, oriented toward its locally more likely
relation. Within each connected component of the crossing graph, the
orientation maximising total satisfied weight is chosen --- equivalently a
minimum-weight feedback arc set. Components up to a configured size are solved
exactly by dynamic programming and larger ones greedily. Relations that cannot
be satisfied are flipped rather than discarded, and both the flip and the
original local probability are recorded, so a reader can see how much the
global correction moved.

\subsection{Layers and metric height}

The corrected relations form a directed acyclic graph whose longest path fixes
the smallest number of layers consistent with them; filaments that never
conflict share a layer, so the layer count follows from the ordering rather
than from the number of instances. Each instance is then placed at a height
respecting its order that keeps filaments from interpenetrating at their
measured radii, with the stack compacted so the film is as thin as those
constraints allow.

\subsection{Configuration}

Every tunable quantity of both stages --- the cost weights and gates above, the
core and window sizes, the two channel weights, the abstention band, and the
exact-solver size limit --- is carried in a single released parameter file,
grouped by the stage it belongs to, with the 2-D grouping section checked
against the frozen record that produced the published numbers.
